% Part: first-order-logic
% Chapter: model-theory
% Section: isomorphism

\documentclass[../../../include/open-logic-section]{subfiles}

\begin{document}

\olfileid{mod}{bas}{iso}
\olsection{సమరూప నిర్మాణాలు}

మొదటిస్థాయి \tetoken{నిర్మాణాలు}{!!{structure}s} రెండు విధాలుగా ఒకేలా ఉండవచ్చు.
ఒక విధంలో అవి ఒకే \tetoken{వాక్యాలను}{!!{sentence}s} సత్యం చేస్తాయి. అటువంటి
\tetoken{నిర్మాణాలను}{!!{structure}s} \emph{మొదటిస్థాయి-వాక్య తుల్యాలు} అంటాం.
కానీ చాలా భిన్నమైన నిర్మాణాలు కూడా ఒకే \tetoken{వాక్యాలను}{!!{sentence}s} సత్యం
చేయవచ్చు---ఉదాహరణకు, వాటిలో ఒకటి \tetoken{లెక్కించదగినది}{!!{enumerable}},
మరొకటి లెక్కించలేనిది కావచ్చు. దీనికి కారణం, ఒక
\tetoken{నిర్మాణపు}{!!a{structure}} అనేక లక్షణాలను మొదటిస్థాయి భాషల్లో వ్యక్తం
చేయలేకపోవడం: భాష తగినంత సమృద్ధిగా లేకపోవచ్చు, లేదా
లొవెన్‌హైమ్--స్కోలెమ్ సిద్ధాంతం వంటి మొదటిస్థాయి తర్కపు మౌలిక
పరిమితులు అడ్డుపడవచ్చు. అందువల్ల \tetoken{నిర్మాణాలు}{!!{structure}s} ఒకేలా
ఉండగల మరొక, మరింత కఠినమైన భావం ఉంది: అవి ఏర్పడిన వస్తువుల్లో మాత్రమే
తేడా ఉండి, నిర్మాణాత్మక లక్షణాల్లో తేడా లేకపోవడం వల్ల అవి మౌలికంగా
ఒకటే కావడం. ఈ భావాన్ని \emph{సమరూపత} ద్వారా ఖచ్చితంగా చెప్పవచ్చు.

\begin{defn}
\ollabel{defn:elem-equiv}
ఒకే \tetoken{భాష}{!!{language}}~$\Lang{L}$కు చెందిన రెండు
\tetoken{నిర్మాణాలు}{!!{structure}s} $\Struct{M}$, $\Struct M'$
ఇవ్వబడ్డాయనుకుందాం. $\Lang{L}$లోని ప్రతి \tetoken{వాక్యం}{!!{sentence}}~$!A$కు
$\Sat{M}{!A}$ అవుతుంది అప్పుడే $\Sat{M'}{!A}$ అయితే, అప్పుడే
$\Struct{M}$ను $\Struct M'$కు \emph{మొదటిస్థాయి-వాక్య తుల్యం} అంటాం;
దీనిని $\Struct{M} \equiv \Struct M'$గా రాస్తాం.
\end{defn}

\begin{defn}
\ollabel{defn:isomorphism}
ఒకే \tetoken{భాష}{!!{language}}~$\Lang L$కు చెందిన రెండు
\tetoken{నిర్మాణాలు}{!!{structure}s} $\Struct{M}$, $\Struct M'$
ఇవ్వబడ్డాయనుకుందాం. కింది షరతులు నెరవేర్చే ప్రమేయం $h \colon
\Domain{M} \to \Domain{M'}$ ఉంటే, అప్పుడే $\Struct{M}$ను
$\Struct M'$కు \emph{సమరూపం} అంటాం; దీనిని $\Struct{M} \simeq
\Struct M'$గా రాస్తాం:
\begin{enumerate}
\item $h$ \tetoken{ఒకటి-ఒకటి}{!!{injective}}: $h(x)=h(y)$ అయితే $x=y$;
\item $h$ \tetoken{సంగ్రస్త}{!!{surjective}}: ప్రతి $y \in \Domain{M'}$కు
  $h(x)=y$ అయ్యే $x \in \Domain{M}$ ఉంటుంది;
\item \ollabel{defn:iso-const}ప్రతి \tetoken{స్థిరసంకేతం}{!!{constant}} $c$కు
  $h(\Assign{c}{M}) = \Assign{c}{M'}$;
\item \ollabel{defn:iso-pred}ప్రతి $n$-స్థాన
  \tetoken{విధేయ సంకేతం}{!!{predicate}}~$P$కు
  \[
  \tuple{a_1, \dots, a_n}\in \Assign{P}{M} \quad\text{అవుతుంది అప్పుడే}\quad
  \tuple{h(a_1), \dots, h(a_n)} \in \Assign{P}{M'};
  \]
\item \ollabel{defn:iso-func}ప్రతి $n$-స్థాన
  \tetoken{ప్రమేయ సంకేతం}{!!{function}} $f$కు
  \[
  h(\Assign{f}{M}(a_1, \dots, a_n)) =
  \Assign{f}{M'}(h(a_1), \dots, h(a_n)).
  \]
\end{enumerate}
\end{defn}

\begin{thm}
\ollabel{thm:isom}
$\Struct{M} \iso \Struct M'$ అయితే, $\Struct{M} \elemequiv
\Struct{M'}$.
\end{thm}

\begin{proof}
$h$ అనేది $\Struct{M}$ నుంచి $\Struct M'$ మీదికి ఒక సమరూపత
అనుకుందాం. ప్రతి కేటాయింపు~$s$కు $h \circ s$ అనేది $h$, $s$ల
సంయోజనం; అంటే $(h \circ s)(x)=h(s(x))$ అయ్యే $\Struct{M'}$లోని
కేటాయింపు. $t$, $!A$లపై ఆగమనంతో కింది మరింత బలమైన వాదనలను
నిరూపించవచ్చు:
\begin{enumerate}
  \item[a.] $h(\Value{t}{M}[s]) = \Value{t}{M'}[h\circ s]$.
  \item[b.] $\Sat{M}{!A}[s]$ అవుతుంది అప్పుడే
    $\Sat{M'}{!A}[h \circ s]$.
\end{enumerate}
మొదటి వాదనను $t$ సంక్లిష్టతపై ఆగమనంతో నిరూపిస్తాం.
\begin{enumerate}
\item $t \ident c$ అయితే, $\Value{c}{M}[s]=\Assign{c}{M}$ మరియు
  $\Value{c}{M'}[h \circ s]=\Assign{c}{M'}$. అందువల్ల
  $h(\Value{t}{M}[s])=h(\Assign{c}{M})=\Assign{c}{M'}$
  (\olref{defn:iso-const} వల్ల; అది \olref{defn:isomorphism}లో ఉంది)
  $=\Value{t}{M'}[h \circ s]$.
\item $t \ident x$ అయితే, $\Value{x}{M}[s]=s(x)$ మరియు
  $\Value{x}{M'}[h \circ s]=h(s(x))$. అందువల్ల
  $h(\Value{x}{M}[s])=h(s(x))=\Value{x}{M'}[h \circ s]$.
\item $t \ident f(t_1, \dots, t_n)$ అయితే,
  \begin{align*}
    \Value{t}{M}[s] &=
      \Assign{f}{M}(\Value{t_1}{M}[s], \dots, \Value{t_n}{M}[s])
      \quad\text{మరియు}\\
    \Value{t}{M'}[h \circ s] &=
      \Assign{f}{M'}(\Value{t_1}{M'}[h \circ s],
      \dots, \Value{t_n}{M'}[h \circ s]).
  \end{align*}
  \sourcecorrection{OLTEMODBAS-003}{రెండవ పంక్తి $\Struct{M'}$లో పదపు విలువను లెక్కిస్తున్నప్పటికీ ఆంగ్ల మూలం అక్కడ $\Assign{f}{M}$ అని రాసింది; సందర్భం, తదుపరి గణనకు అనుగుణంగా దానిని $\Assign{f}{M'}$గా సరిచేశాం.}
  ప్రతి $i$కు $h(\Value{t_i}{M}[s])
  =\Value{t_i}{M'}[h\circ s]$ అనేది ఆగమన పరికల్పన. కాబట్టి
  \begin{align}
    h(\Value{t}{M}[s])
    &= h(\Assign{f}{M}(\Value{t_1}{M}[s], \dots,
    \Value{t_n}{M}[s])) \notag\\
    &= \Assign{f}{M'}(h(\Value{t_1}{M}[s]), \dots,
    h(\Value{t_n}{M}[s])) \ollabel{iso-1}\\
    &= \Assign{f}{M'}(\Value{t_1}{M'}[h \circ s], \dots,
    \Value{t_n}{M'}[h \circ s]) \ollabel{iso-2}\\
    &= \Value{t}{M'}[h\circ s] \notag
  \end{align}
  \sourcecorrection{OLTEMODBAS-004}{ఈ గణన మొదటి పంక్తిలో బయట ఉన్న $h(\,\cdot\,)$కు ఆంగ్ల మూలం ముగింపు కుడి కుండలీకరణను ఇవ్వలేదు; తుల్యతల శ్రేణి సవ్యంగా ఏర్పడేలా ఒక్క \texttt{)}ను పునరుద్ధరించాం.}
  ఇక్కడ \olref{iso-1}, \olref{defn:iso-func} వల్ల; అది
  \olref{defn:isomorphism}లో ఉంది. \olref{iso-2}, ఆగమన పరికల్పన వల్ల
  వస్తుంది.
\end{enumerate}
భాగం (b)ను అభ్యాసంగా వదిలివేస్తున్నాం.

$!A$ ఒక వాక్యమైతే, $s$, $h \circ s$ కేటాయింపులు
అప్రస్తుతం; అప్పుడు $\Sat{M}{!A}$ అవుతుంది అప్పుడే $\Sat{M'}{!A}$.
\end{proof}

\begin{prob}
\olref[mod][bas][iso]{thm:isom}లోని (b) భాగపు నిరూపణను వివరంగా
పూర్తి చేయండి. \olref[mod][bas][iso]{defn:isomorphism}లో సమరూపతలను
లక్షణీకరించే ఐదు ధర్మాల్లో ప్రతి ఒక్కదాన్ని ఎక్కడ వాడారో తప్పక
గమనించండి.
\end{prob}

\begin{defn}
ఒక నిర్మాణం $\Struct{M}$ యొక్క \emph{స్వసమరూపత} అంటే,
$\Struct{M}$ నుంచి దానిమీదికే ఉన్న సమరూపత.
\end{defn}

\begin{prob}
ప్రతి నిర్మాణం $\Struct{M}$కు, $X$ అనేది $\Struct{M}$లో ఎటువంటి
అదనపు ప్రాచలాలు లేకుండా నిర్వచించదగిన ఉపసమితి, $h$ అనేది
$\Struct{M}$ యొక్క స్వసమరూపత
అయితే, $X=\Setabs{h(x)}{x \in X}$ అని (అంటే $h$ కింద $X$ స్థిరంగా
ఉంటుందని) చూపండి.
\sourcecorrection{OLTEMODBAS-010}{ఆంగ్ల అభ్యాసం ``definable`` అని మాత్రమే అంటుంది; ప్రాచలాలతో నిర్వచనాన్ని అనుమతిస్తే వాటిని కదిలించే స్వసమరూపత ఆ నిర్వచిత సమితినీ కదిలించవచ్చు. ప్రతి స్వసమరూపత కింద స్థిరత్వం కలిగించే ఉద్దేశించిన ప్రమేయాన్ని, అదనపు ప్రాచలాలు లేని నిర్వచనానికి స్పష్టంగా పరిమితం చేశాం.}
\end{prob}

\end{document}
